\documentclass[11pt,a4paper]{article}
\usepackage[T1]{fontenc}
\usepackage[utf8]{inputenc}
\usepackage[polish]{babel}
\usepackage{lmodern}
\usepackage[a4paper,margin=2.4cm]{geometry}
\usepackage{hyperref}
\title{HF20.1: naprawa audytu kontraktu i lokalnego MLflow}
\author{Edward Szczypka}
\date{\today}
\begin{document}
\maketitle
\section{Audyt kontraktu}
Komenda CLI musi jawnie importować funkcję
\texttt{audit\_latest\_hourly\_serving\_contract}. Test regresyjny sprawdza
tożsamość obiektu eksportowanego przez CLI i funkcji z modułu audytu.
\section{PowerShell 5.1 i Python}
Wieloliniowy kod Pythona nie jest przekazywany przez \texttt{python -c}.
Zamiast tego uruchamiany jest zwykły plik
\texttt{scripts/enable\_local\_mlflow.py}. Eliminuje to zależność od reguł
cytowania natywnych argumentów Windows PowerShell 5.1.
\section{Local-only}
Walidacja konfiguracji wymusza wyłączenie ObjectStore, uploadu modeli,
publikacji i mirrora manifestu snapshotu w konfiguracji treningowej.
\end{document}
